% GENERATED FILE. Edit canonical lessons, then run npm run generate:books.
% canonical-source-hash: fnv1a64:45f134fad4bd2db4
% canonical-lessons: HI-C500-dulha, HI-C500-mehndi, HI-C500-sagai, HI-C500-barat, HI-C500-vidai

\chapter{A Wedding}
\label{ch:hi-a-wedding}
\hlchaptermodality{500}

\begin{chapteropening}
\textbf{By the end of this chapter:} \emph{I can describe a wedding.}

Complete the last lesson of chapter 500: I can describe a wedding.
\end{chapteropening}

\section[\texorpdfstring{\hi{दूल्हा} (dūlhā)}{दूल्हा (dūlhā)}]{\hi{दूल्हा} (dūlhā) — a bridegroom}

\label{lesson:HI-C500-dulha}



\begin{quote}
Before the new one: say the Hindi for the prime of youth, then the Hindi for to whimper.
\end{quote}

\subsection*{You'll want to know: \hi{दूल्हा}}
\textbf{\hi{दूल्हा}} — \emph{dūlhā} — \textquotedblleft{}a bridegroom\textquotedblright{}.

\begin{quote}
\hi{दूल्हा} \hi{घोड़े} \hi{पर} \hi{बैठकर} \hi{आया।}
\end{quote}

\begin{quote}
\emph{dūlhā ghoṛe par baiṭhkar āyā.}
\end{quote}

\begin{quote}
The bridegroom arrived on horseback.
\end{quote}

Masculine.

\begin{cousinweb}
\textbf{\hi{दूल्हा}} (\emph{dūlhā}) is a tadbhava: CDIAL lists Hindi \emph{dūlhā}, \textquotedblleft{}bridegroom\textquotedblright{}, under Sanskrit \emph{durlabha}, \textquotedblleft{}hard to get, precious\textquotedblright{} (CDIAL 6446): the dear one. The bride is \emph{dulhan}, the feminine of the same word, listed in the same entry.
\end{cousinweb}

\begin{grammarlens}[title={What you've built}]
One more word for describing what you lived and felt.
\end{grammarlens}

\subsection*{Your turn}
\begin{itemize}
\raggedright
  \item \emph{Say it:} \emph{dūlhā}
  \item \emph{Say it:} \emph{dūlhā}, once more
  \item \emph{Say it:} \emph{dūlhā ghoṛe par baiṭhkar āyā.} — The bridegroom arrived on horseback.
  \item \emph{Recall:} say the Hindi for the prime of youth, then the Hindi for to whimper, then say \emph{dūlhā} again
\end{itemize}

\begin{tcolorbox}[breakable,colback=teal!4,colframe=teal!35!black,title={Before you move on}]
What does \hi{दूल्हा} mean? (\textquotedblleft{}A bridegroom\textquotedblright{}.) Say it once more.
\end{tcolorbox}
\section[\texorpdfstring{\hi{मेहंदी} (mehndī)}{मेहंदी (mehndī)}]{\hi{मेहंदी} (mehndī) — henna}

\label{lesson:HI-C500-mehndi}



\begin{quote}
Before the new one: say the Hindi for to whimper, then the Hindi for a bridegroom.
\end{quote}

\subsection*{You'll want to know: \hi{मेहंदी}}
\textbf{\hi{मेहंदी}} — \emph{mehndī} — \textquotedblleft{}henna\textquotedblright{}.

\begin{quote}
\hi{शादी} \hi{के} \hi{दिन} \hi{मीरा} \hi{के} \hi{हाथों} \hi{में} \hi{मेहंदी} \hi{लगी} \hi{थी।}
\end{quote}

\begin{quote}
\emph{śādī ke din Mīrā ke hāthoṁ meṁ mehndī lagī thī.}
\end{quote}

\begin{quote}
On her wedding day Mira's hands were patterned with henna.
\end{quote}

Feminine.

\begin{cousinweb}
\textbf{\hi{मेहंदी}} (\emph{mehndī}) is a tadbhava: CDIAL lists Hindi \emph{mẽhdī}, the henna plant, under Sanskrit \emph{mēndhī} (CDIAL 10328). Its leaves are ground to a paste that leaves red-brown patterns on the hands, above all before a wedding.
\end{cousinweb}

\begin{grammarlens}[title={What you've built}]
One more word for describing what you lived and felt.
\end{grammarlens}

\subsection*{Your turn}
\begin{itemize}
\raggedright
  \item \emph{Say it:} \emph{mehndī}
  \item \emph{Say it:} \emph{mehndī}, once more
  \item \emph{Say it:} \emph{śādī ke din Mīrā ke hāthoṁ meṁ mehndī lagī thī.} — On her wedding day Mira's hands were patterned with henna.
  \item \emph{Recall:} say the Hindi for to whimper, then the Hindi for a bridegroom, then say \emph{mehndī} again
\end{itemize}

\begin{tcolorbox}[breakable,colback=teal!4,colframe=teal!35!black,title={Before you move on}]
What does \hi{मेहंदी} mean? (\textquotedblleft{}Henna\textquotedblright{}.) Say it once more.
\end{tcolorbox}
\section[\texorpdfstring{\hi{सगाई} (sagāī)}{सगाई (sagāī)}]{\hi{सगाई} (sagāī) — an engagement (to marry)}

\label{lesson:HI-C500-sagai}



\begin{quote}
Before the new one: say the Hindi for a bridegroom, then the Hindi for henna.
\end{quote}

\subsection*{You'll want to know: \hi{सगाई}}
\textbf{\hi{सगाई}} — \emph{sagāī} — \textquotedblleft{}an engagement (to marry)\textquotedblright{}.

\begin{quote}
\hi{पिछले} \hi{साल} \hi{मीरा} \hi{की} \hi{सगाई} \hi{हुई।}
\end{quote}

\begin{quote}
\emph{pichle sāl Mīrā kī sagāī huī.}
\end{quote}

\begin{quote}
Mira got engaged last year.
\end{quote}

Feminine.

\begin{cousinweb}
\textbf{\hi{सगाई}} (\emph{sagāī}) is probably the noun of \emph{sagā}, \textquotedblleft{}of full blood relationship\textquotedblright{}, which CDIAL lists under \emph{*svagya-}, a form it reconstructs (CDIAL 13896b): the promise that makes two families kin. CDIAL does not give the \textquotedblleft{}engagement\textquotedblright{} sense itself.
\end{cousinweb}

\begin{grammarlens}[title={What you've built}]
One more word for describing what you lived and felt.
\end{grammarlens}

\subsection*{Your turn}
\begin{itemize}
\raggedright
  \item \emph{Say it:} \emph{sagāī}
  \item \emph{Say it:} \emph{sagāī}, once more
  \item \emph{Say it:} \emph{pichle sāl Mīrā kī sagāī huī.} — Mira got engaged last year.
  \item \emph{Recall:} say the Hindi for a bridegroom, then the Hindi for henna, then say \emph{sagāī} again
\end{itemize}

\begin{tcolorbox}[breakable,colback=teal!4,colframe=teal!35!black,title={Before you move on}]
What does \hi{सगाई} mean? (\textquotedblleft{}An engagement (to marry)\textquotedblright{}.) Say it once more.
\end{tcolorbox}
\section[\texorpdfstring{\hi{बारात} (bārāt)}{बारात (bārāt)}]{\hi{बारात} (bārāt) — the groom's wedding procession}

\label{lesson:HI-C500-barat}



\begin{quote}
Before the new one: say the Hindi for henna, then the Hindi for an engagement (to marry).
\end{quote}

\subsection*{You'll want to know: \hi{बारात}}
\textbf{\hi{बारात}} — \emph{bārāt} — \textquotedblleft{}the groom's wedding procession\textquotedblright{}.

\begin{quote}
\hi{बारात} \hi{में} \hi{सब} \hi{लोग} \hi{नाच} \hi{रहे} \hi{थे।}
\end{quote}

\begin{quote}
\emph{bārāt meṁ sab log nāc rahe the.}
\end{quote}

\begin{quote}
Everyone in the wedding procession was dancing.
\end{quote}

Feminine.

\begin{cousinweb}
\textbf{\hi{बारात}} (\emph{bārāt}) is a tadbhava: CDIAL lists Hindi \emph{barāt}, \textquotedblleft{}marriage procession\textquotedblright{}, under Sanskrit \emph{varayātrā}, built on \emph{vara}, \textquotedblleft{}bridegroom\textquotedblright{}, and \emph{yātrā}, \textquotedblleft{}journey\textquotedblright{} (CDIAL 11322).
\end{cousinweb}

\begin{grammarlens}[title={What you've built}]
One more word for describing what you lived and felt.
\end{grammarlens}

\subsection*{Your turn}
\begin{itemize}
\raggedright
  \item \emph{Say it:} \emph{bārāt}
  \item \emph{Say it:} \emph{bārāt}, once more
  \item \emph{Say it:} \emph{bārāt meṁ sab log nāc rahe the.} — Everyone in the wedding procession was dancing.
  \item \emph{Recall:} say the Hindi for henna, then the Hindi for an engagement (to marry), then say \emph{bārāt} again
\end{itemize}

\begin{tcolorbox}[breakable,colback=teal!4,colframe=teal!35!black,title={Before you move on}]
What does \hi{बारात} mean? (\textquotedblleft{}The groom's wedding procession\textquotedblright{}.) Say it once more.
\end{tcolorbox}
\section[\texorpdfstring{\hi{विदाई} (vidāī)}{विदाई (vidāī)}]{\hi{विदाई} (vidāī) — a farewell; the bride's leave-taking}

\label{lesson:HI-C500-vidai}



\begin{quote}
Before the new one: say the Hindi for an engagement (to marry), then the Hindi for the groom's wedding procession.
\end{quote}

\subsection*{You'll want to know: \hi{विदाई}}
\textbf{\hi{विदाई}} — \emph{vidāī} — \textquotedblleft{}a farewell; the bride's leave-taking\textquotedblright{}.

\begin{quote}
\hi{विदाई} \hi{के} \hi{समय} \hi{मीरा} \hi{और} \hi{उसकी} \hi{माँ} \hi{दोनों} \hi{रोईं।}
\end{quote}

\begin{quote}
\emph{vidāī ke samay Mīrā aur uskī māṁ donõ roīṃ.}
\end{quote}

\begin{quote}
At the bride's farewell, Mira and her mother both wept.
\end{quote}

Feminine.

\begin{cousinweb}
\textbf{\hi{विदाई}} (\emph{vidāī}) is \emph{vidā}, \textquotedblleft{}farewell\textquotedblright{}, with the ending \emph{-ī}; \emph{vidā} is usually traced to Arabic \emph{vidāʿ}, \textquotedblleft{}farewell\textquotedblright{}, which reached Hindi through Persian. At a wedding it is the moment the bride leaves her parents' home.
\end{cousinweb}

\begin{grammarlens}[title={What you've built}]
One more word for describing what you lived and felt.
\end{grammarlens}

\subsection*{Your turn}
\begin{itemize}
\raggedright
  \item \emph{Say it:} \emph{vidāī}
  \item \emph{Say it:} \emph{vidāī}, once more
  \item \emph{Say it:} \emph{vidāī ke samay Mīrā aur uskī māṁ donõ roīṃ.} — At the bride's farewell, Mira and her mother both wept.
  \item \emph{Recall:} say the Hindi for an engagement (to marry), then the Hindi for the groom's wedding procession, then say \emph{vidāī} again
\end{itemize}

\begin{tcolorbox}[breakable,colback=teal!4,colframe=teal!35!black,title={Before you move on}]
What does \hi{विदाई} mean? (\textquotedblleft{}A farewell; the bride's leave-taking\textquotedblright{}.) Say it once more.
\end{tcolorbox}
